%% ma: L12-B08-0122
%% lop: 12
%% chuong: 2
%% bai: 8
%% dang: 12.8.45
%% loai: TN4
%% muc: TH
%% dap_an: "A"
%% nguon: "CĐ2.3.pdf giáo viên gửi 27/10/2026 (M07, Lê Minh Khôi – Nguyễn Anh Khôi, Chương II Chủ đề 3 Biểu thức tọa độ của các phép toán vectơ), Câu 34"
%% kien_thuc: "Bài 8 (khoảng cách hai điểm), lớp 10 (hệ phương trình bậc nhất)"
%% trang_thai: chinh-thuc
%% ghi_chu: "Giáo viên đã duyệt 27/10/2026"
\begin{ex}
Trong không gian $Oxyz$, cho ba điểm $A(2;2;2)$, $B(-2;2;0)$ và $C(4;1;-1)$. Trên mặt phẳng $(Oxz)$, điểm nào dưới đây cách đều ba điểm $A$, $B$, $C$?
\choice{\True $P\left(\dfrac34;0;-\dfrac12\right)$}{$M\left(\dfrac34;0;\dfrac12\right)$}{$Q\left(-\dfrac34;0;\dfrac12\right)$}{$N\left(-\dfrac34;0;-\dfrac12\right)$}
\loigiai{$I(x;0;z)$: $IA^2=IB^2\Leftrightarrow8x+4z=4$; $IA^2=IC^2\Leftrightarrow4x-6z=6$.

Giải hệ: $x=\dfrac34$, $z=-\dfrac12$.}
\end{ex}
